% Kamp III.2, Lemma 9: universal composition of elementary interval conditions.

\subsubsection{Universal composition of elementary interval conditions}
\label{sec:p09-kamp-iii2-lemma9}

Lemma~9 extends the basic universal-family construction to elementary
point conditions and supplied interval representations. The source first
combines each of its three finite disjunctions, then uses Lemma~8 and
proper predicate substitution
\cite[ch.~III, lemma~9, pp.~92--93]{Kamp1968}.

\begin{hbdefinition}{def:p09-kamp-universal-elementary-data}{The three universal families}
Let \(x,y,z\) be pairwise distinct variable names. For finite sizes
\(m,n,p\geq0\), take
\[
\begin{aligned}
 \chi_i&\in\operatorname{E}(z) &&(1\leq i\leq m),\\
 \phi_j&\text{ has property }I(x,z) &&(1\leq j\leq n),\\
 \psi_k&\text{ has property }I(z,y) &&(1\leq k\leq p).
\end{aligned}
\]
Each interval condition has no free variables outside its two indicated
endpoints. A supplied \(I\) representation is part of the construction
data. Let \(\chi,\Phi,\Psi\) be the respective disjunctions,
and define
\[
 U(x,y)\defeq\forall z\,(x<z<y\to\chi(z)\lor\Phi(x,z)\lor\Psi(z,y)).
\]
An empty disjunction is the declared falsehood at an indicated endpoint;
for the unary family it is \(F^K_z\). No empty family deletes the
universal quantifier or either strict inequality. The source's indexed
lists \(0,\ldots,m\), \(0,\ldots,n\), \(0,\ldots,p\)
are its nonempty-family cases; the zero-size convention here is explicit.
\end{hbdefinition}

\begin{hbmethod}{def:p09-kamp-universal-elementary-construction}{Fresh basic words and complete restoration}
Combine all supplied left representations by concatenating their disjunct
lists, and do the same on the right. Obtain the elementary labelled word
of every resulting interval disjunct. The unary disjunction \(\chi\)
has an \(\operatorname{E}(z)\) construction, including its false
empty case.

Choose a fresh unary letter \(Q\) for \(\chi\). At each
occurrence position of every elementary word choose another fresh unary
letter; all chosen letters are distinct and absent from every input
label. Replace each labelled word by the basic past word on those fresh
letters. Apply Lemma~8 to the finite family consisting of \(Q(z)\),
all left basic words at \((x,z)\), and all right basic words at
\((z,y)\).

In its resulting \(I(x,y)\) representation, properly substitute
\(\chi\) for \(Q\), and each original elementary label for its
fresh occurrence letter. Perform these substitutions in a fixed finite
order, with bound variables made fresh before insertion. The replacement
labels contain none of the chosen targets, so later steps do not alter
earlier replacements. Return the resulting finite elementary disjunction.
Every input interval disjunct and every occurrence label is restored.
\end{hbmethod}

\begin{hblemma}{thm:p09-kamp-iii2-lemma9}{proved=annex:p9.kamp.iii2.lemma9}{Kamp's universal elementary-composition lemma}
Under the hypotheses above, \(U(x,y)\) has property \(I(x,y)\).
The stated construction returns actual
\(\operatorname{Ep}(x,y)\) disjuncts whose disjunction is equivalent
to \(U\) under \(x<y\) over \(\mathsf{Com}\).
The claim includes every nonempty family of the source and the explicit
empty-family cases. It uses the given \(I\) equivalences only at
their complete-order scope; no equivalence on incomplete orders is inferred
\cite[ch.~III, lemma~9, pp.~92--93]{Kamp1968}.
\end{hblemma}

\begin{hbexample}{def:p09-kamp-universal-different-disjuncts}{Different disjuncts at different interior points}
Take \(0<1<2<3\), with outer endpoints 0 and 3. Let the unary
condition \(Q(z)\) be false everywhere. Let the left word be
\((A,B,C)\), with \(A=\{0\}\), \(B=T\), \(C=\{1\}\),
and the right word be \((D,E,F)\), with
\(D=\{2\}\), \(E=T\), \(F=\{3\}\), where \(T\)
is this four-point domain. The left disjunct holds at the tested point
1 and fails at 2; the right holds at 2 and fails at 1. Thus the
universal disjunction holds, while neither individual interval disjunct
holds at every interior point. The construction preserves this varying
choice; it does not distribute a universal quantifier over disjunction.

If all three families are empty, the result instead asserts that
\((x,y)\) is empty. On adjacent endpoints it is true; at \((0,3)\)
in this example it is false. The annex verifies both cases.
\end{hbexample}
